% GENERATED FILE. Edit canonical lessons, then run npm run generate:books.
% canonical-source-hash: fnv1a64:806ab230de30d027
% canonical-lessons: PA-C129-hanera, PA-C129-canan, PA-C129-cha, PA-C129-satrangipingh, PA-C129-khargosh

\chapter{Darkness, Light, Shade, Rainbow, Rabbit}
\label{ch:pa-darkness-light-shade-rainbow-rabbit}
\hlchaptermodality{129}

\begin{chapteropening}
\textbf{By the end of this chapter:} \emph{I can say darkness, light, shade, a rainbow and a rabbit.}

Complete the last lesson of chapter 129: I can say darkness, light, shade, a rainbow and a rabbit.
\end{chapteropening}

\section[\texorpdfstring{\pa{ਹਨੇਰਾ} (hanerā)}{ਹਨੇਰਾ (hanerā)}]{\pa{ਹਨੇਰਾ} (hanerā) — darkness}

\label{lesson:PA-C129-hanera}



\begin{quote}
Before the new one: say the Punjabi for the earth, then the Punjabi for moonlight.
\end{quote}

\subsection*{You'll want to know: \pa{ਹਨੇਰਾ}}
\textbf{\pa{ਹਨੇਰਾ}} — \emph{hanerā} — \textquotedblleft{}darkness\textquotedblright{}.

\begin{cousinweb}
\textbf{\pa{ਹਨੇਰਾ}} (\emph{hanerā}) is inherited. CDIAL lists the Punjabi forms \emph{anherā} and \emph{nherā}, \textquotedblleft{}dark; darkness\textquotedblright{}, under \emph{*andhīkāra}, a form beside Sanskrit \emph{andhakāra}, \textquotedblleft{}darkness\textquotedblright{} (CDIAL 386).
\end{cousinweb}

\begin{grammarlens}[title={What you've built}]
One more word for this chapter.
\end{grammarlens}

\subsection*{Your turn}
\begin{itemize}
\raggedright
  \item \emph{Say it:} \emph{hanerā}
  \item \emph{Say it:} \emph{hanerā}, once more
  \item \emph{Say it:} say \emph{hanerā}
  \item \emph{Recall:} say the Punjabi for the earth, then the Punjabi for moonlight, then say \emph{hanerā} again
\end{itemize}

\begin{tcolorbox}[breakable,colback=teal!4,colframe=teal!35!black,title={Before you move on}]
What does \pa{ਹਨੇਰਾ} mean? (\textquotedblleft{}Darkness\textquotedblright{}.) Say it once more.
\end{tcolorbox}
\section[\texorpdfstring{\pa{ਚਾਨਣ} (cānaṇ)}{ਚਾਨਣ (cānaṇ)}]{\pa{ਚਾਨਣ} (cānaṇ) — light}

\label{lesson:PA-C129-canan}



\begin{quote}
Before the new one: say the Punjabi for moonlight, then the Punjabi for darkness.
\end{quote}

\subsection*{You'll want to know: \pa{ਚਾਨਣ}}
\textbf{\pa{ਚਾਨਣ}} — \emph{cānaṇ} — \textquotedblleft{}light\textquotedblright{}.

\begin{cousinweb}
\textbf{\pa{ਚਾਨਣ}} (\emph{cānaṇ}) is a tadbhava, a word inherited and reshaped by sound change. It comes from \emph{*cāndraṇa}, \textquotedblleft{}moonlight\textquotedblright{}, a form CDIAL reconstructs from its descendants (CDIAL 4745). CDIAL lists Punjabi \emph{cānaṇ}, \textquotedblleft{}light\textquotedblright{}.
\end{cousinweb}

\begin{grammarlens}[title={What you've built}]
One more word for this chapter.
\end{grammarlens}

\subsection*{Your turn}
\begin{itemize}
\raggedright
  \item \emph{Say it:} \emph{cānaṇ}
  \item \emph{Say it:} \emph{cānaṇ}, once more
  \item \emph{Say it:} say \emph{cānaṇ}
  \item \emph{Recall:} say the Punjabi for moonlight, then the Punjabi for darkness, then say \emph{cānaṇ} again
\end{itemize}

\begin{tcolorbox}[breakable,colback=teal!4,colframe=teal!35!black,title={Before you move on}]
What does \pa{ਚਾਨਣ} mean? (\textquotedblleft{}Light\textquotedblright{}.) Say it once more.
\end{tcolorbox}
\section[\texorpdfstring{\pa{ਛਾਂ} (chā̃)}{ਛਾਂ (chā̃)}]{\pa{ਛਾਂ} (chā̃) — shade}

\label{lesson:PA-C129-cha}



\begin{quote}
Before the new one: say the Punjabi for darkness, then the Punjabi for light.
\end{quote}

\subsection*{You'll want to know: \pa{ਛਾਂ}}
\textbf{\pa{ਛਾਂ}} — \emph{chā̃} — \textquotedblleft{}shade\textquotedblright{}.

\begin{cousinweb}
\textbf{\pa{ਛਾਂ}} (\emph{chāṁ}) is a tadbhava, a word inherited from Sanskrit and reshaped by sound change. It comes from \emph{chāyā}, \textquotedblleft{}shade, shadow, reflection\textquotedblright{} (CDIAL 5027).
\end{cousinweb}

\begin{grammarlens}[title={What you've built}]
One more word for this chapter.
\end{grammarlens}

\subsection*{Your turn}
\begin{itemize}
\raggedright
  \item \emph{Say it:} \emph{chā̃}
  \item \emph{Say it:} \emph{chā̃}, once more
  \item \emph{Say it:} say \emph{chā̃}
  \item \emph{Recall:} say the Punjabi for darkness, then the Punjabi for light, then say \emph{chā̃} again
\end{itemize}

\begin{tcolorbox}[breakable,colback=teal!4,colframe=teal!35!black,title={Before you move on}]
What does \pa{ਛਾਂ} mean? (\textquotedblleft{}Shade\textquotedblright{}.) Say it once more.
\end{tcolorbox}
\section[\texorpdfstring{\pa{ਸਤਰੰਗੀ} \pa{ਪੀਂਘ} (satrangī pīṅgh)}{ਸਤਰੰਗੀ ਪੀਂਘ (satrangī pīṅgh)}]{\pa{ਸਤਰੰਗੀ} \pa{ਪੀਂਘ} (satrangī pīṅgh) — a rainbow}

\label{lesson:PA-C129-satrangipingh}



\begin{quote}
Before the new one: say the Punjabi for light, then the Punjabi for shade.
\end{quote}

\subsection*{You'll want to know: \pa{ਸਤਰੰਗੀ} \pa{ਪੀਂਘ}}
\textbf{\pa{ਸਤਰੰਗੀ} \pa{ਪੀਂਘ}} — \emph{satrangī pīṅgh} — \textquotedblleft{}a rainbow\textquotedblright{}.

\begin{cousinweb}
\textbf{\pa{ਸਤਰੰਗੀ} \pa{ਪੀਂਘ}} (\emph{satrangī pīṅgh}) is \textquotedblleft{}the seven-coloured swing\textquotedblright{}. \emph{Sat}, \textquotedblleft{}seven\textquotedblright{}, comes from Sanskrit \emph{sapta} (CDIAL 13139); \emph{raṅgī} is built on \emph{raṅg}, \textquotedblleft{}colour\textquotedblright{}; and \emph{pīṁgh}, \textquotedblleft{}swing\textquotedblright{}, comes from Sanskrit \emph{preṅkha}, \textquotedblleft{}swing\textquotedblright{} (CDIAL 8996).
\end{cousinweb}

\begin{grammarlens}[title={What you've built}]
One more word for this chapter.
\end{grammarlens}

\subsection*{Your turn}
\begin{itemize}
\raggedright
  \item \emph{Say it:} \emph{satrangī pīṅgh}
  \item \emph{Say it:} \emph{satrangī pīṅgh}, once more
  \item \emph{Say it:} say \emph{satrangī pīṅgh}
  \item \emph{Recall:} say the Punjabi for light, then the Punjabi for shade, then say \emph{satrangī pīṅgh} again
\end{itemize}

\begin{tcolorbox}[breakable,colback=teal!4,colframe=teal!35!black,title={Before you move on}]
What does \pa{ਸਤਰੰਗੀ} \pa{ਪੀਂਘ} mean? (\textquotedblleft{}A rainbow\textquotedblright{}.) Say it once more.
\end{tcolorbox}
\section[\texorpdfstring{\pa{ਖ਼ਰਗੋਸ਼} (khargosh)}{ਖ਼ਰਗੋਸ਼ (khargosh)}]{\pa{ਖ਼ਰਗੋਸ਼} (khargosh) — a rabbit}

\label{lesson:PA-C129-khargosh}



\begin{quote}
Before the new one: say the Punjabi for shade, then the Punjabi for a rainbow.
\end{quote}

\subsection*{You'll want to know: \pa{ਖ਼ਰਗੋਸ਼}}
\textbf{\pa{ਖ਼ਰਗੋਸ਼}} — \emph{khargosh} — \textquotedblleft{}a rabbit\textquotedblright{}.

\begin{cousinweb}
\textbf{\pa{ਖ਼ਰਗੋਸ਼}} (\emph{khargosh}) is a loan, usually traced to Persian \emph{khargōsh}, \textquotedblleft{}hare\textquotedblright{}. It is built from \emph{khar}, \textquotedblleft{}donkey\textquotedblright{}, and \emph{gōsh}, \textquotedblleft{}ear\textquotedblright{}.
\end{cousinweb}

\begin{grammarlens}[title={What you've built}]
One more word for this chapter.
\end{grammarlens}

\subsection*{Your turn}
\begin{itemize}
\raggedright
  \item \emph{Say it:} \emph{khargosh}
  \item \emph{Say it:} \emph{khargosh}, once more
  \item \emph{Say it:} say \emph{khargosh}
  \item \emph{Recall:} say the Punjabi for shade, then the Punjabi for a rainbow, then say \emph{khargosh} again
\end{itemize}

\begin{tcolorbox}[breakable,colback=teal!4,colframe=teal!35!black,title={Before you move on}]
What does \pa{ਖ਼ਰਗੋਸ਼} mean? (\textquotedblleft{}A rabbit\textquotedblright{}.) Say it once more.
\end{tcolorbox}
